\chapter{软件工程}

\section{领域导言与演进背景}

\begin{preamblebox}[知识领域概述与时代演进]
早在 20 世纪 70 年代初，软件工程先驱戴夫·帕纳斯（David L. Parnas）便提出了振聋发聩的经典定义：\textbf{“软件工程是多人协同构建多版本程序的过程”}。这一洞见揭示了软件工程的核心本质：如果说程序设计是决定我们编写单个程序能力的技能，那么软件工程则在\textbf{“时间”}与\textbf{“人员”}两个决定性维度上与单纯的编程产生了根本分野。

\textbf{维度一：人员协同}\\
软件工程项目本质上是一项团队协作事业。仅仅成为一名独来独往的编程高手是远远不够的。合格的软件工程师必须具备良好的沟通、倾听、同理心与协作能力。编程或许可以是一项高度个体化的智力活动，但软件工程从立项之初就是一项社会化协作活动，与专业操守、团队运作、跨职能交流和同侪评审紧密交织。

\textbf{维度二：跨越时间的生命周期（Multi-Version \& Lifespan）}\\
软件工程项目通常是“多版本演进”的。软件拥有其预期生命周期，往往需要在数月、数年甚至数十年中持续稳定运转。在此期间，为了满足不断变化的产品需求，功能特性会被频繁增删；底层的计算基础设施、操作系统、编译器与编程语言版本会不断升级演变，外部第三方依赖库亦会频繁迭代。与课堂上“写完即弃”的单人作业相比，这种直面时间流逝与环境变迁的现实是软件工程特有的挑战：\textbf{仅仅构建一个当前能跑通的程序是远远不够的，系统必须能够经受住时间的考验，在漫长的演化中始终保持正确运转}。当今技术界最棘手的一系列工程挑战，其共同根源皆在于“时间必然催生变化”：向后兼容性、版本偏差/漂移、依赖地狱与依赖管理、数据库模式演变以及通信协议演进。

\textbf{软件工程高等教育的特有挑战：}\\
在象牙塔的学术环境中开展软件工程教学面临着独特的结构性难题。鉴于软件工程与传统编程的核心差异在于“时间”与“团队”，在有限的学期周数内设计出既能充分展现团队协作摩擦、又能真实还原长周期演进风险的教学实验极其困难。此外，软件工程中的许多高阶主题，只有当学习者在真实工业场景——无论是前沿企业实习、深度参与大型开源社区、还是全职工程岗位——中亲身体验过大规模协作与长期代码维护后，才能产生最深刻的共鸣与领悟。哪怕仅一个月的全职实战经验，其有效工程工时往往就已超过一门普通大学课程的全学期学时。

因此，CS2023 软件工程课程规范明确聚焦于：\textbf{绝大多数计算机专业应届毕业生走向工业界首日即迫切需要的核心理念}。这些内容要么对于仅受过传统单人编程训练的学生而言是全新盲区，要么是工业界日常协作中习以为常、却极少被显式系统传授的抽象工程法则。涵盖的核心主题包括：
\begin{itemize}
    \item 工业级软件测试与质量验证
    \item 团队协作与冲突消解
    \item 专业化工程沟通与代码评审
    \item 软件架构、API 契约与详细设计
    \item 软件维护、重构与长周期演化
    \item 现代自动化工程工具链
\end{itemize}

\textbf{教学模式权衡：讲授与项目驱动的辩证融合}\\
部分理论内容非常适合采用标准课堂讲授或“讲授+实验”的形式进行剖析。深入探讨版本控制系统的底层理论模型（如 Git 的有向无环图 DAG）、工业界主流分支演进策略、静态代码分析与动态插桩分析的理论边界、以及依赖拓扑网络中的“菱形依赖问题”，能够帮助学生建立扎实的技术心智模型。

另一方面，软件工程的大量核心精髓必须在亲手实践中才能内化。长期以来，大作业项目驱动型课程是此类教学的通用载体。然而，CS2023 特别指出其中的潜在教学风险：\textbf{如果学生唯一的协作经验是与其他未经受过工程训练的新手组成团队去完成一个单一项目，他们极易对团队协作的真正复杂性与最佳实践产生认知偏差与错误归因}。因此，任课教师必须在理论沉淀与项目实践之间审慎权衡。规范强烈建议项目课程倡导“敏捷迭代与极速反馈”——重复演练若干小步快跑的规范任务（例如遵循敏捷模式构建），其教学效果远优于让新手在高度摩擦中狼狈应对一个庞大单体项目。若学校不具备开设长周期综合项目的条件，任何能够让学生体验协作与时间维度的轻量级设计都极具价值：例如向大型既有代码库贡献特性、不同小组分别开发系统的异构模块并通过严格契约对接、结对分别编写编解码器进行联调互测等。

\textbf{时代背景与团队工程的必然性：}\\
所有现实证据表明，软件在整个人类社会运行中的中枢地位仍将持续高速膨胀，而“车库里两个极客黑客敲出一套伟大软件”的英雄主义神话时代早已宣告终结。当今工业界所有重要的软件系统无一例外都是高度组织化团队协同的产物，绝大多数开发工作都是在既有庞大代码资产之上进行演进与复用。软件工程教学是计算机专业学生从象牙塔走向现实工程世界的决定性桥梁——我们必须让每位学子深刻认识到：在工业实践中，几乎没有任何现实软件是作为单打独斗的个人项目从零编写的。沟通、协作、规划、测试与工具链，才是支撑他们走出校园并在广阔现实世界中立足的真正核心基石。

\textbf{自 CS2013 以来的核心变革：}
\begin{enumerate}
    \item \textbf{高度聚焦现代工业软件工作流}：大幅强化版本控制、自动化测试套件、规范化同侪代码评审以及现代 CI/CD 自动化工具链的必修要求。
    \item \textbf{弱化传统宏观项目管理与团队领导力}：大幅削减以往偏向高级项目经理与团队主管的抽象管理说教（如大型组织架构规划、重型软件过程度量），使课程精准匹配应届毕业生的首要岗位诉求。
    \item \textbf{聚焦团队参与、专业沟通与协同编码}：将核心落脚点全面转向作为团队普通成员如何高效参与敏捷协作、开展无指责技术沟通、规范处理代码冲突与接口演进。
\end{enumerate}
\end{preamblebox}

\section{核心学时分配建议}

本知识领域在“向日葵模型”下的核心学时推荐如下表所示（计算机核心共 6 学时，领域核心共 21 学时，合计 27 核心学时）：

\begin{table}[H]
\centering
\setlength{\tabcolsep}{3.5pt}\fontsize{8pt}{9.6pt}\selectfont\begin{tabularx}{\columnwidth}{Y c c}
\toprule
\textbf{知识单元} & \textbf{计算机核心} & \textbf{领域核心} \\
\midrule
团队协作与沟通 & 2 & 2 \\
工具与开发环境 & 1 & 3 \\
产品需求工程 & -- & 2 \\
软件架构与系统设计 & 1 & 4 \\
软件构造与工程实现 & 1 & 3 \\
软件验证与确认 & 1 & 3 \\
代码重构与软件演进 & -- & 2 \\
软件可靠性工程 & -- & 2 \\
形式化方法 & -- & \emph{拓展选修} \\
\midrule
\textbf{合计} & \textbf{6 学时} & \textbf{21 学时} \\
\bottomrule
\end{tabularx}
\end{table}

\section{知识单元与详细大纲}

\subsection{团队协作与沟通}
\textbf{【学时建议】}：2 计算机核心学时，2 领域核心学时。

由于传统初级编程教学的个体化特征，绝大多数初修软件工程的学生极少或从未接触过软件工程的协作本质。单纯的实战练习（如课程大作业）虽有助益，但必须在课堂讲授与研讨中系统注入关于清晰、高效、专业沟通与团队协作核心价值的深度剖析。

\subsubsection{计算机核心主题}
\begin{itemize}
    \item \textbf{高效团队沟通机制}：同步沟通与异步沟通的权衡、技术备忘录撰写、沟通中的倾听与建设性反馈机制。
    \item \textbf{团队冲突根因与消解策略}：技术分歧与人际摩擦的本质辨析、对事不对人的专业态度（针对技术提案和代码实现提出批评，坚决避免人身攻击）、团队冲突仲裁与共识达成路径。
    \item \textbf{协同编程工程实践}：
    \begin{itemize}
        \item 结对编程：驾驶员与领航员角色轮换、心智协同与即时审查。
        \item 同侪代码评审：评审礼仪、建设性评注、评审清单、小粒度提交与及时响应。
    \end{itemize}
    \item \textbf{软件团队中的角色与职责划分}：
    \begin{itemize}
        \item 团队协同的规模放大优势与互补效应。
        \item 协作带来的沟通成本、同步阻碍与组织复杂性（人月神话与布鲁克斯定律直觉）。
    \end{itemize}
    \item \textbf{团队开发流程与敏捷过程规范}：
    \begin{itemize}
        \item 任务认领与责任追踪。
        \item 敏捷工作量估算与故事点。
        \item 高效站会与会议结构（Meeting Structure、议程管理、时间盒机制）。
        \item 工作排期与迭代节奏。
    \end{itemize}
    \item \textbf{团队多元化的战略价值}：认知背景多元、技能互补、视角包容度对软件设计鲁棒性与创新的巨大推动。
\end{itemize}

\subsubsection{领域核心主题}
\begin{itemize}
    \item \textbf{作为团队整体与外部利益相关者高效对接}：
    \begin{itemize}
        \item 管理层及其他非技术业务团队（产品、运营、法务、市场营销）。
        \item 商业客户（客户需求协商、合同边界澄清）。
        \item 最终用户（真实使用体验反馈、可用性痛点挖掘）。
    \end{itemize}
    \item \textbf{物理集中、全分布式、混合制与虚拟远程团队的系统性风险}：
    \begin{itemize}
        \item 跨时区与异步通信断层、认知脱节与孤岛效应。
        \item 组织架构脆弱性、单点人员故障与技术断层。
        \item 远程协作风险的缓解策略、应急恢复预案与团队心理健康维护。
    \end{itemize}
\end{itemize}

\subsubsection{示范性学习成果与能力目标}
\begin{itemize}
    \item 严格遵循高效团队专业沟通规范进行技术交流与工作协同。
    \item 阐明团队冲突的根源、潜在危害与建设性价值——特别聚焦于能够在不攻击他人人格的前提下，就技术思想或架构方案开展富有成效的良性争鸣。
    \item 在真实的团队工程环境中主导或推动行之有效的冲突化解策略。
    \item 在协同开发与结对编程过程中展现高水准的协作效能与默契度。
    \item 在软件开发团队中合理提议、协商并分配关键角色与工程职责。
    \item 针对团队技术评审或敏捷例会，撰写结构严谨的会议议程并在会议中严格执行。
    \item 深入参与团队综合项目，全面践行团队建设核心要素、建立健康包容的团队工程文化，并制定与推进团队工作执行计划。
    \item 深刻阐释并积极推动多元化与包容性为软件研发团队带来的深远工程益处与创新潜能。
    \item 阐明软件团队在技术与非技术双重层面与外部利益相关者建立定期沟通机制的重要性与核心策略。
    \item 系统梳理集中式、全分布式、混合制与虚拟远程团队所面临的特有工程风险与潜在故障点，并制定相应的风险缓释、容错与故障复盘机制。
\end{itemize}

\subsection{现代软件工程工具与开发环境}
\textbf{【学时建议】}：1 计算机核心学时，3 领域核心学时。

在过去的一代人时间里，整个软件工业界对软件工程工具链的依赖呈现爆发式增长：分布式版本控制系统已无处不在，自动化测试框架成为标配，静态代码分析与动态插桩检测融入日常，持续集成（CI）系统更是成为工程交付的生命线。与此相伴，日益强大的集成开发环境（IDE）提供了毫秒级代码语义检索、全局索引、自动化安全重构以及对各类外部工程工具的无缝集成。深刻理解这些工业级工具的运行机制与选型权衡——尤其是版本控制系统——是现代软件工程师不可或缺的底层素养。

\subsubsection{计算机核心主题}
\begin{itemize}
    \item \textbf{软件配置管理（SCM）与版本控制基础}：
    \begin{itemize}
        \item 代码与配置纳管：在版本控制中统一管理配置、可复现构建与环境即代码。
        \item 分支演进策略：功能开发分支vs 生产发布分支；工业级“基于主干的开发”范式。
        \item 代码分支集成策略：合并与变基的底层原理、适用场景与提交历史整洁性权衡。
    \end{itemize}
\end{itemize}

\subsubsection{领域核心主题}
\begin{itemize}
    \item \textbf{发布管理}：发布流水线、灰度发布、版本打标与发布分支维护。
    \item \textbf{专业测试工具链}：单元测试执行器、测试框架、静态代码分析器与动态内存/并发分析工具。
    \item \textbf{软件研发全流程自动化}：
    \begin{itemize}
        \item 构建系统：高速增量构建、密封隔离构建（Hermetic Builds，消除宿主机隐式污染）、构建缓存、现代主流构建工具体系对比（Make、Gradle、Bazel 等）。
        \item 持续集成 (CI)：利用自动化构建流水线与自动化测试套件作为工程“单一事实来源”，即时验证主干最新代码的构建可达性与基础测试通过率。
        \item 依赖管理：上游与第三方依赖库自动解析、依赖冲突消解、包管理器生态、语义化版本控制（SemVer，主版本号.次版本号.修订号）的严格契约。
    \end{itemize}
    \item \textbf{设计与协作沟通工具}：架构设计文档系统、图表绘制工具、通用建模语言规范（如 UML 类图、时序图、架构拓扑图）。
    \item \textbf{现代工具集成概念与机制}：Webhook、自动化钩子脚本、流水线编排协议与工具间上下文打通。
    \item \textbf{现代 IDE 高级工程套件}：全项目符号索引与毫秒级搜索、交互式断点调试、自动化代码重构、代码片段模板与辅助工具深度集成。
\end{itemize}

\subsubsection{示范性学习成果与能力目标}
\begin{itemize}
    \item 深入辨析集中式配置管理系统（如 SVN）与分布式版本控制系统（如 Git）在架构拓扑、协作效率与容灾能力上的核心差异。
    \item 阐述如何综合运用分支管理、标签与版本控制策略支持工业级软件发布管理。
    \item 在小型团队协作项目中准确识别关键配置项，并熟练使用分布式源代码版本控制工具完成协同开发。
    \item 阐明主流静态代码检查工具与动态测试工具如何无缝集成至开发环境与日常编码流水线中。
    \item 深刻理解持续集成（CI）系统作为团队共享主干代码健康状态唯一基准事实源的战略地位。
    \item 全面阐述在为特定软件系统进行工具链选型时的核心考量要素，涵盖需求追踪、架构建模、编码实现、构建自动化与自动化测试等各个环节。
    \item 展现熟练运用现代化软件工具链支持中等规模软件产品全流程开发的工程实践能力。
\end{itemize}

\subsection{产品需求工程}
\textbf{【学时建议】}：2 领域核心学时。

知道“如何构建系统”固然重要，但如果我们不知道“该构建什么”，一切技术投入都将付诸东流。产品需求（亦称需求工程、产品设计、需求诱导、产品需求规格说明 PRD 等）系统引导学生学习围绕定义新产品或新功能特性所必需的完整规范化流程。

\subsubsection{领域核心主题}
\begin{itemize}
    \item \textbf{功能需求的结构化与场景化表述}：
    \begin{itemize}
        \item 运用用例或用户故事（User Stories，如“作为……我希望……以便于……”格式）精确描述功能需求。
        \item 掌握至少一种规范化记录与组织功能需求的方法体系。
        \item 深刻理解需求如何作为后续架构设计、编码实现与验收测试的直接输入。
        \item 辩证评估特定需求描述方法的表达力优势与潜在局限。
    \end{itemize}
    \item \textbf{需求规约的核心质量属性}：一致性、合法有效性、完备性与工程可行性。
    \item \textbf{需求诱导与获取}：
    \begin{itemize}
        \item 需求来源辨析：最终用户、系统管理员、运维人员、合规审查方与业务支持人员。
        \item 需求采集方法：问卷调查、深度访谈、焦点小组、现场观察与埋点用户行为大数据分析。
    \end{itemize}
    \item \textbf{非功能需求与质量属性}：系统安全性、可用性/易用性、高并发性能、可靠性、可维护性与合规隐私要求（交叉 SEC 知识领域）。
    \item \textbf{需求阶段的风险识别与风险管理}：需求模糊度评估、范围蔓延防御与技术可行性预警。
    \item \textbf{需求规格说明书的专业沟通与形式化}：半形式化规范、验收准则与契约化签署。
\end{itemize}

\subsubsection{拓展选修主题}
\begin{itemize}
    \item \textbf{敏捷原型构建}：低保真线框图、高保真交互原型与最小可行产品（MVP），作为需求启发、推演与用户确认的核心工具。
    \item \textbf{产品需求演进与变更控制}：当业务需求发生不可逆变更时，系统化评估波及范围、依赖链路与需修改的系统模块。
    \item \textbf{敏捷工程估算}：
    \begin{itemize}
        \item 掌握科学评估开发任务工作量的经典技术（斐波那契故事点、计划扑克 Planning Poker、专家 Delphi 法）。
        \item 在实训中持续开展工作量预估，并与实际消耗工时进行复盘校准。
        \item 认识到软件估算的极高难度——即使学生初期存在显著偏差，亲历估算偏差并在复盘中修正认知模型亦极具教学价值。
    \end{itemize}
\end{itemize}

\subsubsection{示范性学习成果与能力目标}
\begin{itemize}
    \item 沿采集成本、覆盖广度、主观偏差等多个正交维度系统对比不同的需求获取方法。
    \item 深入辨析两种功能需求描述方法（如传统用例规范 vs 敏捷用户故事/客户访谈）的本质差异及其各自适配的工程场景。
    \item 给定一组明确的产品需求规范与一组候选系统行为，准确裁定哪些系统行为是被强制要求的、被允许的或被严格禁止的。
    \item 针对一个简易软件系统，独立完成全套业务需求的采集、分类与系统化整理。
    \item 给定现有系统的架构描述与一组新增业务需求，准确识别出系统中必须发生适应性调整的核心模块与代码范围。
    \item 从纷繁杂乱的需求集合中，精准辨识并分离功能需求与非功能需求（质量属性）。
    \item 针对待建软件系统快速构建原型（Mock-up 原型或 MVP），以实证方式验证关键业务需求集的可行性。
    \item 对一组研发任务预估工时，并在项目实施后将初始预估值与实际消耗工时进行量化复盘对比。
    \item 依据任务间的逻辑依赖关系制定合理的任务实现时序路线图，以“尽早消除核心技术风险”为第一准则组织开发。
    \item 为简易软件系统撰写专业、规范且具备可验证性的软件需求规格说明书（SRS/PRD）。
\end{itemize}

\subsection{软件架构与系统设计}
\textbf{【学时建议】}：1 计算机核心学时，4 领域核心学时。

产品需求聚焦于软件系统面向最终用户的“外部功能”，而软件设计则聚焦于面向软件工程师的“内部组件拓扑与结构逻辑”。本单元全面涵盖宏观的软件体系结构设计以及微观层面的 API 接口与数据契约设计。

\subsubsection{计算机核心主题}
\begin{itemize}
    \item \textbf{系统设计核心原则}：
    \begin{itemize}
        \item 抽象层次划分：宏观架构设计与微观详细设计。
        \item 关注点分离：业务逻辑、数据持久化与用户展现界面的解耦。
        \item 信息隐藏：隐藏内部数据表示与算法细节，仅暴露稳定抽象。
        \item 耦合度与内聚度：追求“高内聚、低耦合”的工程本质与量化直觉。
    \end{itemize}
    \item \textbf{经典软件体系结构与设计范式}：
    \begin{itemize}
        \item 核心设计范式：自顶向下功能分解/分层架构、面向数据架构、面向对象分析与设计（OOAD）、事件驱动架构。
        \item 经典标准体系结构：客户端/服务器（C/S）、多层企业架构、管道-过滤器、模型-视图-控制器（MVC）。
        \item 组件边界界定与依赖拓扑图谱识别。
    \end{itemize}
    \item \textbf{大规模编程与小规模编程的本质分野}：Programming in the Large（聚焦模块交互、契约制定、团队协作与边界治理）vs. Programming in the Small（聚焦单体算法逻辑、局部循环与数据结构操纵）。
    \item \textbf{代码异味与内在质量表征}：辨析代码“功能正确性”与“内在工程质量/可维护性”的本质区别（过长函数、上帝类、霰弹式修改、重复代码等）。
\end{itemize}

\subsubsection{领域核心主题}
\begin{itemize}
    \item \textbf{工业级 API 设计核心准则}：
    \begin{itemize}
        \item 一致性：接口命名、参数顺序与错误语义在系统内部严格统一，符合领域通用惯用法（降低学习成本，减少误用概率）。
        \item 可组合性：小而美的正交接口可通过流水线或组合轻松构建复杂能力。
        \item 契约式设计与文档规范：清晰界定接口对系统状态的变更影响而非泄露其内部实现细节；显式声明前置条件、后置条件与类不变量。
        \item 可扩展性：在不破坏现有调用的前提下为未来能力拓展预留平滑空间。
        \item 健壮的错误报告机制：错误语义明确、可预测且具备可操作性；非法入参必须显式报错拒绝；能够内部安全自愈的偶发故障无需过度上报。
    \end{itemize}
    \item \textbf{识别并形式化数据不变量与时间/状态不变量}：状态机转换边界、并发序列时序约束。
    \item \textbf{软件设计的结构模型与行为模型}：类拓扑图、调用关系网、状态迁移图。
    \item \textbf{数据架构设计}：核心数据结构选型、持久化存储引擎选型与交互协议（交叉 DM 知识领域）。
    \item \textbf{需求可追溯性}：建立业务需求项与架构设计组件之间的双向追溯矩阵。
\end{itemize}

\subsubsection{拓展选修主题}
\begin{itemize}
    \item \textbf{高级设计建模方法}：统一建模语言（UML）类图、实体-联系图（E-R 图）、时序图与组件部署图。
    \item \textbf{软件设计质量的量化度量与技术债务评估}：圈复杂度、模块不稳定性、抽象度指标。
    \item \textbf{安全设计与编码原则（Secure Design Principles，交叉 SEC 领域）}：
    \begin{itemize}
        \item 最小特权原则。
        \item 故障安全默认原则。
        \item 心理可接受度原则。
    \end{itemize}
    \item \textbf{系统设计多维权衡决策}：执行效率 vs. 可靠性、极致安全 vs. 操作便捷性等维度的科学决断。
\end{itemize}

\subsubsection{示范性学习成果与能力目标}
\begin{itemize}
    \item 给定一个高层架构方案，准确辨析其采用的标准软件体系结构模式（如 MVC、分层、微内核、事件驱动）。
    \item 针对简易软件系统选择并运用适宜的设计范式开展系统设计，并清晰阐明关注点分离、信息隐藏等设计原则在其中的贯彻落实。
    \item 识别并重构存在严重缺陷的系统设计方案，消除违反关注点分离与信息隐藏的设计弊端。
    \item 在给定的宏观架构设计中，精确识别所有软件组件之间的依赖拓扑与调用链条。
    \item 为大型软件系统中的独立组件设计生产级 API，包括形式化定义并归档其所有操作的不变量、调用契约与异常错误条件。
    \item 从一致性、可组合性与可扩展性等维度，严谨评估一份 API 设计规约的工程优劣。
    \item 对现有系统设计进行平滑扩展，在不破坏既有架构约束的前提下接纳全新业务功能模块。
    \item 设计一套兼顾访问效率与内存开销的高质量内部数据结构，以完美实现给定的公共 API 接口表面。
    \item 审核给定的软件架构设计，精确判定其完全满足了哪些原始业务需求项。
    \item 将自然语言描述的业务系统设计规约准确转化为规范的 UML 类图与交互时序图。
    \item 调整并加固存在安全缺陷的架构设计，使其严格遵循最小特权原则与故障安全默认原则。
    \item 从执行性能、资源消耗、可靠性与易用性等多元质量维度，系统对比与评估两种不同的软件设计方案。
\end{itemize}

\subsection{软件构造与工程实现}
\textbf{【学时建议】}：1 计算机核心学时，3 领域核心学时。

软件构造聚焦于直接指导高质量代码产出的工业实践：实用的测试驱动、规范的代码文档、严谨的编程规范、异常处理与防御性实现。进阶主题涵盖安全编码、依赖注入、任务时序与生产级既有代码库演进。

\subsubsection{计算机核心主题}
\begin{itemize}
    \item \textbf{实用小粒度单元测试实践}：
    \begin{itemize}
        \item 单元测试：针对独立函数或单一类型编写具备高分支覆盖率的自动化测试用例。
        \item 测试驱动开发 (TDD)：TDD 对学生具有不可替代的心理学教育价值——在编写实现之前，预先挑战自己去探寻目标 API 的刁钻边界、异常输入与极端测试用例，其心智阻力显著更小。如果先写完实现，开发者往往会本能地避开破坏自己心血作品的极端输入；而“测试先行”则让学生在发现潜在问题中获得智力满足，并在开发过程中亲眼目睹测试用例从红变绿，获得极大的工程成就感。
    \end{itemize}
    \item \textbf{工业级工程文档编写规范}：
    \begin{itemize}
        \item 接口契约文档：准确描述接口的功能约束、前后置条件与数据不变量。
        \item 实现级文档规范：重点聚焦于棘手、微妙且非直观的代码逻辑（无论是因为运用了高级语言特性还是业务逻辑本身高度复杂）。坚决杜绝在显而易见的简单语法或通用操作上添加无效注释；注释应致力于阐明数据流向与设计意图（代码在做什么、为什么这样做）；对于晦涩代码应优先通过重构使其自解释，必要时辅以精准注释点拨核心玄机。
    \end{itemize}
\end{itemize}

\subsubsection{领域核心主题}
\begin{itemize}
    \item \textbf{工业级编码风格与规范}：
    \begin{itemize}
        \item 团队代码规范（Style Guides，如 Google/PEP8 风格指南）。
        \item 规范化注释、代码格式化工具统一约束。
        \item 严谨自解释的标识符命名规范。
    \end{itemize}
    \item \textbf{高质量编码“最佳工程实践”}（交叉 安全隐私/防御性编程与 SDF 知识领域）：
    \begin{itemize}
        \item 防御性编程实践：输入合法性断言、边界严格截断、空值防御。
        \item 安全编码核心原则：抵御注入攻击、权限提升防御。
        \item 运用现代异常处理机制提升软件系统的健壮性与容错能力。
    \end{itemize}
    \item \textbf{代码调试技术}：符号调试器、条件断点、核心转储分析与调用栈溯源。
    \item \textbf{第三方成熟库与开源框架的工程化集成}：依赖评估、许可证合规、封装适配。
\end{itemize}

\subsubsection{拓展选修主题}
\begin{itemize}
    \item \textbf{更大规模的工程级测试技术}：
    \begin{itemize}
        \item 测试替身：桩代码、模拟对象、伪实现的本质差异与选型。
        \item 依赖注入 (DI)：通过控制反转解耦组件依赖，大幅提升单元测试的可测性。
    \end{itemize}
    \item \textbf{研发任务时序编排与里程碑规划}：
    \begin{itemize}
        \item 任务依赖识别：准确梳理复杂工程中不同任务间的上下游强依赖拓扑。
        \item 里程碑设定：设定具有明确功能闭环与交付价值的进度检查点。
        \item 风险前置消减：敏锐识别项目中风险最高的未知领域，优先安排排雷任务。
    \end{itemize}
    \item \textbf{典型底层安全漏洞防御}：缓冲区与整型溢出、数据竞争与并发竞态条件、特权滥用与未安全初始化、外部不可信输入校验。
    \item \textbf{自动化 API 文档生成技术}（如 Doxygen, Javadoc, Sphinx）。
    \item \textbf{开发工程上下文辨析}：全新“绿地开发”vs. 工业级“既有遗留代码库演进”——变更波及影响分析与渐进式变更落地。
    \item \textbf{软件发布管理工程流}。
\end{itemize}

\subsubsection{示范性学习成果与能力目标}
\begin{itemize}
    \item 针对小型软件组件（包含若干函数或单一实体类），编写完备、规范的自动化单元测试用例。
    \item 为小型软件组件撰写规范详尽的对外接口文档，并在代码内部关键复杂逻辑处编写精准的实现注释。
    \item 阐述为实现高可靠性、高执行效率与鲁棒容错设计所必需的技术、编码习惯用法与底层机制。
    \item 熟练运用语言级异常处理机制，编写具备弹性恢复能力的鲁棒代码。
    \item 详细阐述工业级安全编码规范与防御性编程的核心指导原则。
    \item 在小型软件开发项目中选择并严格贯彻执行统一的代码风格规范。
    \item 系统对比与辨析组件集成策略，包括自顶向下、自底向上以及三明治（混合）集成模式的优劣。
    \item 阐述在针对特定项目全新开发的定制代码库中，进行代码变更分析与规范化落地的全流程。
    \item 阐述针对工业界庞大复杂的既有遗留代码库进行代码变更分析与安全实施的工程规范。
    \item 重构一段存在隐患的简易程序，彻底根除其中的常见安全漏洞（如缓冲区溢出、整数溢出和竞态条件）。
    \item 编写一个能够执行非平凡复杂任务、且对外部非法输入和运行时偶发异常具备极高容忍弹性的软件组件。
\end{itemize}

\subsection{软件验证与确认 (软件测试)}
\textbf{【学时建议】}：1 计算机核心学时，3 领域核心学时。

软件验证与确认（V\&V）聚焦于如何最大化发挥测试的工程价值——深刻理解测试的战略定位、系统的失效模式、以及优秀测试与劣质测试之间的本质差异。

\subsubsection{计算机核心主题}
\begin{itemize}
    \item \textbf{验证与确认的哲学辨析}：
    \begin{itemize}
        \item 验证：“Are we building the thing right?”（我们是否以正确的方式构建了系统？系统是否符合预设规格说明？）。
        \item 确认：“Did we build the right thing?”（我们是否构建了正确的系统？系统是否真正解决了用户的实际需求？）。
    \end{itemize}
    \item \textbf{测试为何至关重要}：软件组件在漫长周期的持续代码演进中能否始终保持功能正常？测试是重构与演化的唯一安全气囊。
    \item \textbf{多维测试目标体系}：系统易用性、运行可靠性、规约一致性、系统吞吐与响应性能、全栈安全性（Security，交叉 SEC 领域）。
    \item \textbf{工业级测试层级划分}：单元测试、集成测试、确认/验收测试、端到端系统测试。
    \item \textbf{测试代码与生产代码的代码风格哲学差异}：
    \begin{itemize}
        \item DAMP（Descriptive And Meaningful Phrases，富有描述性与自解释性）vs. DRY（Don't Repeat Yourself，不要重复自己）。
        \item 在测试代码中，教条地消除一切重复反而会破坏用例的自包含性与可读性；为了使失败用例的上下文一目了然，适度的结构性代码重复在测试工程中是完全正当且值得提倡的。
    \end{itemize}
\end{itemize}

\subsubsection{领域核心主题}
\begin{itemize}
    \item \textbf{测试规划与用例设计生成}：
    \begin{itemize}
        \item 基于形式化模型与需求规约生成严密的测试用例。
        \item 多维测试覆盖度：测试矩阵；代码覆盖率（语句、分支、路径覆盖）；环境覆盖度（跨硬件 CPU 架构、异构 OS、主流浏览器与运行时版本）。
        \item 测试数据设计与边界输入生成（等价类划分、边界值分析）。
    \end{itemize}
    \item \textbf{现代测试开发技术体系}（交叉 软件开发基础/开发方法与 软件工程/软件构造）：
    \begin{itemize}
        \item 测试驱动开发（TDD）深度实践。
        \item 面向对象测试、测试打桩与依赖注入。
        \item 黑盒测试与白盒测试技术。
        \item 现代化测试工具集成：代码覆盖率探针、静态代码扫描、模糊测试（Fuzzing）。
    \end{itemize}
    \item \textbf{研发全生命周期中的 V\&V 闭环}：
    \begin{itemize}
        \item 同侪代码评审。
        \item 测试自动化与工具链全自动调度。
        \item 提交前门禁测试vs. 提交后持续集成测试。
        \item 测试深度覆盖率与测试执行吞吐量/排队延迟之间的工程权衡。
        \item 缺陷追踪、严重级别划分与缺陷重现机制。
    \end{itemize}
    \item \textbf{特定前沿领域的 V\&V 独特挑战}：高并发与性能基准测试、异步与分布式并行系统、高可靠安全攸关系统、高精度数值计算验证。
\end{itemize}

\subsubsection{拓展选修主题}
\begin{itemize}
    \item \textbf{高级 V\&V 工具与自动化进阶}：静态分析引擎、模糊测试（Fuzzing）、动态内存/竞态检测工具、缺陷日志自动化分析。
    \item \textbf{测试终止准则与缺陷密度估算}：基于缺陷播种的剩余故障估算、伪随机数在测试生成中的科学运用。
    \item \textbf{性能基准评测与剖析技术}：
    \begin{itemize}
        \item 吞吐量与延迟量化。
        \item 极限负载下的降级与压力测试（FIFO 与 LIFO 请求调度策略）。
        \item 加速比与强/弱可扩展性计算（阿姆达尔定律 Amdahl's Law、古斯塔夫森定律 Gustafson's Law）。
        \item 性能评估关键指标确立。
        \item 常见性能瓶颈诊断：计算受限、内存带宽受限、I/O 与网络延迟受限。
        \item 基准测试的统计学最佳实践：测量不确定度估算、置信区间、多轮采样消除噪声抖动、高精度硬件时钟计时挑战。
        \item 性能测试数据的专业可视化图表呈现。
    \end{itemize}
    \item \textbf{复杂异步、并行与分布式并发系统的测试方案}。
    \item \textbf{非代码工程工件（技术文档、培训手册、合规凭据）的验证与确认}。
\end{itemize}

\subsubsection{示范性学习成果与能力目标}
\begin{itemize}
    \item 深刻阐释测试在软件生命周期与长期演进中的核心战略价值。
    \item 严格辨析软件“验证”与“确认”的本质哲学与工程差异。
    \item 详细阐述软件测试涵盖的多维核心目标（可用性、可靠性、性能、安全等）。
    \item 系统对比不同类型与层级的测试策略（回归测试、单元测试、集成测试、系统测试与用户验收测试）。
    \item 描述制定完备测试计划与高效生成测试用例的核心工程方法。
    \item 针对一段中等规模代码，制定涵盖测试矩阵、边界测试数据与异常输入的完整测试计划。
    \item 落地执行该测试计划，编写并运行自动化测试套件。
    \item 阐明测试驱动开发（TDD）的核心原则，并解释自动化测试在 TDD 循环中扮演的核心支柱角色。
    \item 深入剖析面向对象软件测试所面临的特有挑战（多态分派、隐式状态、对象生命周期）。
    \item 阐述模拟测试与依赖注入（DI）的机制及其在隔离外部依赖中的具体应用。
    \item 作为团队协作成员，对同事提交的中等规模代码段执行专业严谨的同侪代码评审。
    \item 阐释现代化工具链在软件质量确认与合规性证明中所能发挥的关键作用。
    \item 为小型软件项目搭建全自动化测试执行与报表流水线。
    \item 深入分析提交前测试与提交后测试各自的定位、收益与算力开销权衡。
    \item 辩证研讨测试用例覆盖率深度与测试流水线执行吞吐量/反馈时延之间的工程权衡及其对敏捷开发的影响。
    \item 熟练运用工业级缺陷追踪系统在小型项目中闭环管理软件缺陷的生命周期。
    \item 探讨传统经验测试在并发、分布式与安全攸关等特定前沿领域中所暴露的理论与实践局限。
    \item 详细阐释并横向对比多种高级验证确认工具的技术原理与适用场景。
    \item 在小型软件项目中实现多种静态扫描与动态测试工具的自动化集成执行。
    \item 说明在测试用例生成中如何以及何时科学使用随机数，并规避测试不可复现风险。
    \item 阐明软件潜在残留故障估算的核心数学理论与工程方法。
    \item 基于缺陷播种法与缺陷密度历史数据，科学估算小型软件应用中的剩余缺陷数量。
    \item 清晰定义吞吐量与响应时延，并针对具体系统架构各列举典型实例。
    \item 解释计算加速比与不同形式的扩展性（强扩展与弱扩展）的计算模型与理论边界。
    \item 准确诊断并归纳常见的计算密集型、内存受限型与 I/O 延迟型性能瓶颈。
    \item 详述构建工业级性能基准测试时的统计学方法准则与实验设计规范。
    \item 阐明在高精度性能测量中高分辨率计时的技术方案与面临的系统抖动挑战。
    \item 为小型项目确立关键性能评价指标，构建并运行基准测试，完成统计学显著性分析并绘制专业可视化报表。
    \item 阐明测试异步、并发与并行软件系统的特有挑战与调试应对策略。
    \item 为包含异步或并发代码的中等规模模块编制包含测试矩阵与竞争诱发数据的测试计划。
    \item 阐述对非代码软件工程工件（如用户手册、架构白皮书）进行系统化验证与确认的评审流程。
\end{itemize}

\subsection{代码重构与软件演进}
\textbf{【学时建议】}：2 领域核心学时。

软件系统的长期维护与演化是软件工程区别于单纯编程的核心特征。本单元聚焦于系统演进中的核心规律、代码重构技艺、向后兼容性契约与工业级版本控制规范。

\subsubsection{领域核心主题}
\begin{itemize}
    \item \textbf{海勒姆定律（Hyrum's Law / 隐式接口定律）}：
    \begin{quote}
    \emph{“当一个 API 拥有足够多的使用者时，系统在契约中承诺的内容已不再重要：你所展现出的所有可被观察的行为（即使未在文档中声明、甚至是偶然出现的内部未定义行为），都会被某些下游用户所依赖。”}
    \end{quote}
    深刻理解这一定律对于大型系统接口设计、内部行为修改以及变更风险评估具有决定性指导意义。
    \item \textbf{向后兼容性的本质探析}：
    \begin{itemize}
        \item 兼容性并非某单一组件的内在单体属性，而是一种双方契约关系的相对属性。
        \item 向后兼容性必须在“提供者+ 消费者”的具体业务场景中，或者基于提供方明确承诺并达成的兼容性模型下进行科学评估。
    \end{itemize}
    \item \textbf{代码重构工程实践}：
    \begin{itemize}
        \item 经典重构模式：重命名、函数内联、函数提炼、抽取类/接口等。
        \item 现代 IDE 自动化重构工具的安全应用与原理。
        \item 静态分析工具在自动检测代码坏味道、提出重构建议以及批量生成安全补丁中的实战应用。
    \end{itemize}
    \item \textbf{版本控制与演进规范}：
    \begin{itemize}
        \item 语义化版本控制：Major.Minor.Patch 的严格契约与破坏性变更声明。
        \item 基于主干的持续演进与短生命周期特性分支。
    \end{itemize}
\end{itemize}

\subsubsection{拓展选修主题}
\begin{itemize}
    \item \textbf{大规模软件重构}：
    \begin{itemize}
        \item 当重构改动规模过于庞大无法原子提交（跨千万行代码库），或无法在提供者与所有消费者之间原子同步变更时（跨多代码仓库、下游存在私有闭源消费者）的工业级渐进式重构技术。
        \item 双轨并存机制：同时保留新旧两套 API 使之平滑共存。
        \item 微步演进：将庞大重构拆解为系列微小的行为保持变换。
        \item 区分“可直接感知的 API 消费者”与“不可见/隐匿的外部消费者（Consumers I can't see）”带来的不同发布安全策略。
    \end{itemize}
\end{itemize}

\subsubsection{示范性学习成果与能力目标}
\begin{itemize}
    \item 准确识别特定接口的显式契约行为与所有隐式可观察行为，并依据海勒姆定律深刻评估破坏隐式行为可能带来的严重工程崩溃风险。
    \item 综合静态代码分析工具的告警报告与软件设计准则，敏锐定位代码库中急需重构劣质代码。
    \item 严谨研判某项代码修改是否属于真正的“向后兼容”，并能给出关于哪些具体调用场景受支持、哪些不受支持的清晰声明。
    \item 在对既有下游调用方产生零影响或极低波及影响的前提下，重构接口的内部实现以显著改善代码清晰度、内聚度与可维护性。
    \item 结合当前项目所采用的版本规范策略（如 SemVer），严密评估某项提议的系统变更是否具备足够的发布安全性。
    \item 制定周密的多阶段平滑迁移方案，在超大规模代码库中安全地修改、替换并最终废弃某核心 API 的旧有默认行为。
\end{itemize}

\subsection{软件可靠性工程}
\textbf{【学时建议】}：2 领域核心学时。

软件系统的可靠性关乎现代社会基础设施的安全平稳运转。本单元系统剖析可靠性的数学本质、全生命周期缺陷防御、系统容错设计、以及现代站点可靠性工程（SRE）的先进理念。

\subsubsection{领域核心主题}
\begin{itemize}
    \item \textbf{可靠性核心概念体系}：失效率概率、平均无故障时间（MTBF）、平均修复时间（MTTR）；故障、错误与失效的演化因果链路。
    \item \textbf{面向不同领域软件的可靠性需求定级}（交叉 SEP 知识领域：消费级应用 vs. 医疗/航空/金融级关键系统）。
    \item \textbf{软件失效的缺陷根源与防御闭环}：软件失效根本上由代码缺陷诱发，因此高可靠性的目标是追求生产环境缺陷最少化——首要在源头减少缺陷注入（通过专业培训、系统化架构规划、高质量编码规范），其次在交付前最大化消除注入的缺陷（严密测试、静态扫描、同侪评审）。
    \item \textbf{软件可靠性、系统级可靠性与失效行为特征}（交叉 系统基础/基于冗余实现系统可靠性）。
    \item \textbf{缺陷注入与消除全周期}：各阶段缺陷检出率度量与消除手段选型。
    \item \textbf{“错误预算”与“零缺陷”理念的深度思辨}：SRE 现代错误预算理念（允许合理可控的服务中断以换取极速创新迭代）与传统绝对零缺陷理念的对比，准确划定二者各自的适用产业领域。
\end{itemize}

\subsubsection{拓展选修主题}
\begin{itemize}
    \item \textbf{软件可靠性增长模型}：基于失效历史数据的可靠性预测曲线。
    \item \textbf{软件容错架构与容错模型}：
    \begin{itemize}
        \item 容错上下文的深刻工程差异：在飞行控制等关键系统中，必须竭尽全力避免系统崩溃，通过冗余算路自愈；而在数据流批处理系统中，在损坏的脏数据落盘之前让系统主动崩溃、快速失败反而是最具价值的安全容错机制。
    \end{itemize}
    \item \textbf{软件可靠性工程（SRE）工业级实践}：混沌工程、故障注入演练、轻量级模型检验。
    \item \textbf{强相关与相互独立失效域识别}及其对分布式系统整体可用性的连锁影响。
    \item \textbf{基于度量与可观测性的可靠性分析}：全链路遥测、指标监控与智能告警、实时 SLO/SLA 仪表盘、版本发布准入量化指标。
\end{itemize}

\subsubsection{示范性学习成果与能力目标}
\begin{itemize}
    \item 阐述如何为特定业务场景下的软件系统科学测定与协商其所必需的可靠性级别（如服务等级目标 SLO）。
    \item 深刻阐明在真实工程中试图实现极高可靠性指标（如“五个九”99.999\% 可用性）所面临的指数级成本与物理阻碍。
    \item 掌握在软件全生命周期的每个关键阶段（需求、架构、编码、测试、交付、运维）最小化缺陷注入与最大化缺陷清除的系统化工程方法。
    \item 展现综合运用多种可靠性预测模型对特定软件系统的可靠性演进趋势进行量化评估的能力。
    \item 识别并设计能够达成特定可靠性级别指标的软件容错体系结构。
    \item 熟练运用多种维度的冗余机制（时间冗余、空间冗余、数据冗余、逻辑异构冗余）实现软件容错。
    \item 在系统架构设计中敏锐识别单点故障 (SPOF)依赖并提出根除方案。
\end{itemize}

\subsection{形式化方法 (拓展选修)}
\textbf{【学时建议】}：本单元内容为拓展选修，高度依赖离散结构领域的基础数学知识（特别是 DS/基础逻辑与 DS/证明技术）。

形式化方法运用数理逻辑与形式语言对软件规约与行为进行严格的数学证明与系统化模型分析，是保障关键系统万无一失的高级工程手段。

\subsubsection{核心主题}
\begin{itemize}
    \item \textbf{接口形式化规约}：
    \begin{itemize}
        \item 前置条件、后置条件与状态不变量的形式化定义。
        \item 编写与分析形式化前后置条件的专用规范语言与证明工具系统。
    \end{itemize}
    \item \textbf{形式化方法高度适用的核心工程问题域}：
    \begin{itemize}
        \item 无锁并发编程、微观数据竞争隐患。
        \item 复杂异步与分布式共识系统（死锁 Deadlock、活锁 Livelock、饥饿、状态空间爆炸）。
    \end{itemize}
    \item \textbf{形式化分析与经验缺陷检测手段的多维对比}：与传统单元测试、集成测试、模糊测试（Fuzzing）在覆盖完备性与工程投入产出比上的核心权衡。
    \item \textbf{现代软件形式化建模与自动化分析工具}：
    \begin{itemize}
        \item 模型检验器（Model Checkers，如 TLA+、SPIN）。
        \item 模型发现器（Model Finders，如 Alloy）。
    \end{itemize}
\end{itemize}

\subsubsection{示范性学习成果与能力目标}
\begin{itemize}
    \item 阐释形式化规约与自动化分析技术在大型复杂安全攸关软件系统研发中的关键角色，系统对比其与经验测试在验证和确认层面的本质优势与局限。
    \item 将形式化规约与形式化验证工具切实应用于低复杂度的软件核心设计、并发算法或关键业务协议之中。
    \item 深入剖析在工业工程实践中引入形式化规范语言所能带来的极高可靠性收益，及其伴随的高昂学习门槛与建模维护成本。
\end{itemize}

\section{专业素养与思维倾向}

软件工程的学科特质决定了从业者不仅需要掌握扎实的计算机科学技术，更必须具备特定的职业素养、思维方式与心理倾向：

\begin{itemize}
    \item \textbf{协作精神}：软件工程正日益被公认为一项高强度的“团队运动”——合格的软件工程师必须能够与他人高效协同。\textbf{谦逊、尊重与信任}是构筑健康团队协作关系的基石，也是在软件工程领域取得长远成功的根本前提。
    \item \textbf{专业操守}：软件工程师构建的数字化系统深刻重塑着全球数十亿人的日常生活与经济活动。清醒意识到软件系统在现代社会中的重大影响力、恪守崇高的工程伦理准则，并在团队内部与外部日常跨部门交互中始终展现专业、体面与受人尊重的行为风范，是从业者的立身之本。
    \item \textbf{敏锐沟通}：在庞大复杂的现代软件工程项目中，没有任何单一工程师能够通晓系统的所有技术细节。成功的软件项目无不依赖工程师之间高频、清晰、坦诚与严谨的沟通，以实现高保真度的工程协调与敏捷推进。
    \item \textbf{严谨细致}：软件工程对细节的高度敏感性要求团队中的每一位成员都展现出高度一致的严谨态度。在该领域的卓著成就无一例外得益于一丝不苟的工程习惯——对系统规约的全面理解、对规范化流程的严格遵循、以及对投机取巧、偷工减料的本能抵制。
    \item \textbf{责任担当}：软件工程的高度协作性进一步放大了个人担当的核心价值。推诿扯皮、不兑现交付承诺、对阻塞性问题隐瞒不报，是导致团队摩擦升级、技术债务激增乃至软件项目彻底流产的经典祸根。
\end{itemize}

\section{数学基础要求}

\textbf{建议先修内容：}
\begin{itemize}
    \item \textbf{应用概率统计基础}：用于系统多轮性能测试数据对比、A/B 测试实验评估、用户调研问卷结果统计推断、以及性能基准测试中的测量不确定度与置信区间分析。
\end{itemize}

\section{交叉主题与共享概念}

\subsection{共享概念}
\begin{itemize}
    \item \textbf{职业素养与团队协作}：与社会、伦理与专业素养（SEP）知识领域深度共享。
    \item \textbf{数据建模}：与数据管理（DM）知识领域深度共享。
    \item \textbf{安全编程与防御性开发}：与计算机安全（SEC）知识领域深度共享。
\end{itemize}

\subsection{跨学科贯穿主题}
\begin{itemize}
    \item \textbf{工程伦理与社会责任}：深刻贯穿于需求分析、合规审查、错误预算制定与系统可靠性权衡的全生命周期中。
    \item \textbf{工业级编程实现}：作为一切软件架构设计落地与单元测试演练的代码级载体。
\end{itemize}

\section{课程体系规划与开课建议}

针对高校计算机科学专业高年级开设的《高级软件工程》专业核心课程，规范建议至少包含以下 30 学时的进阶教学模块：
\begin{itemize}
    \item 团队协作与沟通 —— 4 学时
    \item 工具与开发环境 —— 4 学时
    \item 产品需求工程 —— 2 学时
    \item 软件架构与系统设计 —— 5 学时
    \item 软件构造与工程实现 —— 4 学时
    \item 代码重构与软件演进 —— 2 学时
    \item 软件可靠性工程 —— 2 学时
    \item 融合来自 SEP 领域的“工程伦理与专业沟通”模块 —— 7 学时
\end{itemize}

\textbf{先修课程要求}：
\begin{itemize}
    \item \textbf{软件开发基础 (SDF)}：学生必须具备足够扎实的个人代码编写流利度，唯有如此，他们才能将心智精力从繁琐的基础语法调试中解放出来，全力聚焦于团队协作、规范流程、架构权衡与工程演进。
\end{itemize}

\textbf{核心能力目标陈述}：\\
成功修完本门课程的合格毕业生，应当能够熟练为团队同事开展高质量的同侪代码评审（尤其注重专业友善的沟通礼仪与团队心理安全感建设），能够独立编写结构完备的自动化单元测试套件，娴熟驾驭现代软件工程工具链（工业级 IDE、分布式版本控制系统、静态分析与覆盖率检查工具），并全面胜任软件工程团队对应届初级工程师所期望的一切核心专业工程交付活动。

\section{真实能力任务规范}

\begin{taskbox}[任务 1：研发中新增业务代码的自动化测试套件构建]
\textbf{【任务情境】}：在敏捷迭代中为正在开发的新核心业务功能编写自动化测试，确保新代码在当前及未来长期维护中始终保持功能正确。\\
\textbf{【能力陈述】}：剖析目标 API 的预期行为与设计契约，将其系统分解为涵盖典型输入、极端边界输入与非法异常输入的测试用例集，构建完整的自动化单元测试套件，确保该功能在后续长期代码演进中持续保持正确性与稳健性。\\
\textbf{【能力领域】}：软件 / 应用 \quad \textbf{【能力单元】}：需求分析 / 系统设计 / 编码实现 / 质量测试 / 技术文档\\
\textbf{【知识单元关联】}：软件开发基础/开发方法、软件工程/软件构造 \quad \textbf{【达标等级】}：\textbf{应用与开发}
\end{taskbox}

\begin{taskbox}[任务 2：为团队同事执行规范化同侪代码评审]
\textbf{【任务情境】}：团队成员提交了一项核心功能实现的代码拉取请求（PR），需要你在合并入主干前开展严格的技术把关。\\
\textbf{【能力陈述】}：秉持专业素养与同理心，以清晰、建设性且富有协作精神的方式与提交代码的同事沟通，就代码的逻辑正确性、测试完备性、潜在安全隐患、代码异味与架构可读性提供详尽切实的评审反馈。\\
\textbf{【能力领域】}：软件 / 应用 \quad \textbf{【能力单元】}：系统设计 / 技术文档 / 代码改进\\
\textbf{【知识单元关联】}：软件开发基础/开发方法、软件工程/团队协作、软件工程/软件验证与确认 \quad \textbf{【达标等级】}：\textbf{应用}
\end{taskbox}

\begin{taskbox}[任务 3：工业级 API 接口文档编制与契约形式化]
\textbf{【任务情境】}：为系统关键业务模块提取通用服务接口，供团队内外其他研发团队进行跨模块调用集成。\\
\textbf{【能力陈述】}：准确梳理目标组件的 API 表面接口，编写结构严谨、表意精准的对外技术文档，明确界定参数类型、前置条件、后置条件、系统状态副作用与异常错误码规范，确保接口被正确理解与安全调用。\\
\textbf{【能力领域】}：软件工程 \quad \textbf{【能力单元】}：需求分析 / 系统设计 / 编码实现 / 技术文档\\
\textbf{【知识单元关联】}：软件开发基础/基础编程概念、软件工程/软件构造、软件工程/软件设计 \quad \textbf{【达标等级】}：\textbf{应用}
\end{taskbox}

\begin{taskbox}[任务 4：遗留与第三方代码库的缺陷排查与调试修复]
\textbf{【任务情境】}：线上运行的遗留业务系统出现复杂异常崩溃，该模块由前人编写且缺少完整注释，需紧急定位并修复。\\
\textbf{【能力陈述】}：认识到在工业界“阅读理解代码的频率远高于编写全新代码”的现实规律，能够熟练通读并透彻理解他人编写的陌生代码，运用现代调试工具链与调用栈追踪技术进行严密审视排查，准确定位隐蔽 Bug 根源并安全重构修复。\\
\textbf{【能力领域】}：软件 / 应用 \quad \textbf{【能力单元】}：系统实现 / 评估诊断 / 软件维护 / 代码改进\\
\textbf{【知识单元关联】}：软件开发基础/开发方法、软件工程/软件构造 \quad \textbf{【达标等级】}：\textbf{应用}
\end{taskbox}

\begin{taskbox}[任务 5：针对研发、设计与调试场景合理进行现代工具选型]
\textbf{【任务情境】}：团队启动一项全新复杂分布式服务的研发，需要搭建兼顾研发效能、代码质量与协作透明度的工具链平台。\\
\textbf{【能力陈述】}：全面洞察现代主流软件工程工具族（静态分析器、动态插桩检测器、分布式版本控制、代码覆盖率分析器、自动化重构工具等）的原理与特性，针对具体的工程痛点与开发瓶颈，科学评估并选用最适宜的工具组合解决问题。\\
\textbf{【能力领域】}：软件 / 应用 \quad \textbf{【能力单元】}：系统实现 / 工具集成\\
\textbf{【知识单元关联】}：软件开发基础/开发方法、软件工程/工具与环境、软件工程/软件构造、软件工程/软件验证与确认 \quad \textbf{【达标等级】}：\textbf{解释}
\end{taskbox}

\begin{taskbox}[任务 6：生产环境核心 API 的平滑演进与安全重构]
\textbf{【任务情境】}：生产环境中已被海量下游系统广泛依赖的基础核心 API 需要调整底层架构并变更部分默认行为。\\
\textbf{【能力陈述】}：在充分考虑海勒姆定律的前提下，精准评估语法变更与语义行为变动的潜在破坏风险，制定周密的多阶段迁移与双轨平滑演化方案，运用现代化重构工具与语义版本机制安全落地重构，向利益相关者清晰传达迁移指引。\\
\textbf{【能力领域】}：软件 / 应用 \quad \textbf{【能力单元】}：系统设计 / 软件维护\\
\textbf{【知识单元关联】}：软件工程/代码重构与软件演进 \quad \textbf{【达标等级】}：\textbf{应用}
\end{taskbox}

\begin{taskbox}[任务 7：自动化代码工程格式化与规范统一约束]
\textbf{【任务情境】}：在大型团队多人协同开发中，解决不同成员编码格式不一引发的代码冲突与 Review 噪音。\\
\textbf{【能力陈述】}：熟练掌握行业通用代码规范与自动化代码格式化工具，在团队流水线与本地 Git Hooks 中落地统一格式约束，深刻践行工程一致性与整洁代码的核心价值。\\
\textbf{【能力领域】}：软件 / 应用 \quad \textbf{【能力单元】}：系统实现 / 软件维护\\
\textbf{【知识单元关联】}：软件开发基础/基础编程概念、软件工程/软件构造 \quad \textbf{【达标等级】}：\textbf{应用}
\end{taskbox}

\begin{taskbox}[任务 8：复杂业务组件的高内聚低耦合 API 架构设计]
\textbf{【任务情境】}：根据给定的复杂业务需求说明，设计一套对外交互的高内聚、低耦合工业级 API。\\
\textbf{【能力陈述】}：综合权衡多维质量属性（执行性能、扩展性、安全性、易用性），设计一套逻辑清晰、命名一致、具备高度可组合性的公共 API 接口集，并给出形式化数据不变量与契约保障。\\
\textbf{【能力领域】}：软件工程 \quad \textbf{【能力单元】}：需求分析 / 架构设计\\
\textbf{【知识单元关联】}：软件工程/软件设计 \quad \textbf{【达标等级】}：\textbf{开发}
\end{taskbox}

\begin{taskbox}[任务 9：关键系统架构设计决策与权衡技术文档撰写]
\textbf{【任务情境】}：面对技术选型存在重大分歧的架构设计，撰写一份正式的架构决策记录供全员评审。\\
\textbf{【能力陈述】}：深刻理解软件设计本质上是针对多种质量属性（运行性能、可用性、维护成本、代码可读性等）的权衡艺术，撰写严谨规范的技术文档，系统阐释重大设计决策背后的技术考量、放弃的备选方案与根本决策依据。\\
\textbf{【能力领域】}：软件 / 应用 \quad \textbf{【能力单元】}：需求分析 / 架构设计 / 技术文档\\
\textbf{【知识单元关联】}：软件工程/软件设计 \quad \textbf{【达标等级】}：\textbf{开发}
\end{taskbox}

\begin{taskbox}[任务 10：作为高效成员参与敏捷团队协作与全周期协同]
\textbf{【任务情境】}：作为初级研发工程师加入敏捷开发团队，参与长周期的产品迭代开发。\\
\textbf{【能力陈述】}：成为团队中高产且可靠的团队成员，展现优秀的团队沟通技巧与长期心理建设意识，主动分担任务，积极参与敏捷站会与复盘，在团队面临工期压力与摩擦时维持良性人际互动与高质量工程交付。\\
\textbf{【能力领域】}：工程应用 \quad \textbf{【能力单元】}：技术文档 / 团队社会议题调适 / 过程改进\\
\textbf{【知识单元关联】}：软件工程/团队协作 \quad \textbf{【达标等级】}：\textbf{应用}
\end{taskbox}

\begin{taskbox}[任务 11：面向开放式综合工程项目的跨角色协同设计与敏捷交付]
\textbf{【任务情境】}：跨专业组队承接来自非技术外部客户提出的模糊开放式需求项目，需在限定预算与工期内完成交付。\\
\textbf{【能力陈述】}：与团队同侪及非技术客户紧密协同，梳理并澄清模糊需求，在严格的时间表和有限资源约束下，完成端到端的系统架构设计、编码实现、持续集成与部署验证，成功达成既定项目目标。\\
\textbf{【能力领域】}：软件 / 应用 \quad \textbf{【能力单元】}：需求工程 / 架构设计 / 技术文档\\
\textbf{【知识单元关联】}：软件工程/团队协作、软件工程/软件设计、软件工程/软件构造 \quad \textbf{【达标等级】}：\textbf{应用}
\end{taskbox}

\section{工作委员会成员名单}

\textbf{委员会主席}：
\begin{itemize}
    \item Titus Winters, 谷歌公司 —— \emph{注：《Google软件工程：源自世界顶级代码库的实践经验》资深领衔作者}
\end{itemize}

\textbf{委员会委员}：
\begin{itemize}
    \item Brett A. Becker, 都柏林大学
    \item Adam Vartanian, Cord 公司
    \item Bryce Adelstein Lelbach, 英伟达公司
    \item Patrick Servello, CIWRO 气象雷达研究所
    \item Pankaj Jalote, 德里因陀罗普拉斯塔信息技术学院
    \item Christian Servin, 埃尔帕索社区学院
\end{itemize}

\textbf{学术与工业界特别贡献者}：
\begin{itemize}
    \item Hyrum Wright, 谷歌公司 —— \emph{注：软件工程著名“海勒姆定律（Hyrum's Law）”命名提出者}
    \item Olivier Giroux, 苹果公司
    \item Gennadiy Civil, 谷歌公司
\end{itemize}

\clearpage
\onecolumn
\section{附录：核心主题与技能要求等级对照表}

\begin{longtable}{p{0.20\textwidth} p{0.46\textwidth} c c c}
\toprule
\textbf{知识单元} & \textbf{核心知识点与具体主题} & \textbf{认知等级} & \textbf{分类} & \textbf{学时} \\
\midrule
\endfirsthead
\toprule
\textbf{知识单元} & \textbf{核心知识点与具体主题} & \textbf{认知等级} & \textbf{分类} & \textbf{学时} \\
\midrule
\endhead
\bottomrule
\endfoot
\multirow{2}{*}{\makecell[l]{软件工程/团队协作\\与沟通}} 
& 有效沟通、团队冲突根因与消解策略、结对编程、同侪代码评审、角色分工、敏捷流程、团队多元化 & 评估 & 计算机核心 & 2 \\
\cmidrule{2-5}
& 与外部利益相关者（管理层/客户/用户）对接、物理与远程虚拟团队风险剖析 & 解释 & 领域核心 & 2 \\
\midrule
\multirow{2}{*}{\makecell[l]{软件工程/工具与\\开发环境}} 
& 软件配置管理与版本控制、可复现构建、基于主干开发的分支策略、合并与变基 & 评估 & 计算机核心 & 1 \\
\cmidrule{2-5}
& 发布管理、静态/动态测试工具、构建系统、CI 持续集成流水线、SemVer 依赖管理、现代 IDE 高级能力 & 解释 & 领域核心 & 3 \\
\midrule
\makecell[l]{软件工程/产品需求工程} 
& 用例与用户故事、需求一致性/完备性/可行性、需求诱导采集、非功能质量属性、需求风险管理与规约形式化 & 应用 & 领域核心 & 2 \\
\midrule
\multirow{2}{*}{\makecell[l]{软件工程/软件架构\\与系统设计}} 
& 系统设计核心原则（抽象、关注点分离、信息隐藏、高内聚低耦合）、经典架构模式（C/S、分层、MVC）、大/小规模编程辨析、代码异味 & 解释 & 计算机核心 & 1 \\
\cmidrule{2-5}
& API 设计原则（一致性、可组合性、契约设计、错误报告）、数据与状态不变量、结构与行为模型、数据架构设计、需求可追溯性 & 应用 & 领域核心 & 4 \\
\midrule
\multirow{2}{*}{\makecell[l]{软件工程/软件构造\\与工程实现}} 
& 实用小型测试（单元测试、测试驱动开发 TDD 心理学机制）、规范化接口文档与微妙实现注释 & 应用 & 计算机核心 & 1 \\
\cmidrule{2-5}
& 工业级编码风格与规范指南、高质量最佳实践（防御性编程、安全编码、异常容错）、断点调试、第三方成熟库集成 & 应用 & 领域核心 & 3 \\
\midrule
\multirow{2}{*}{\makecell[l]{软件工程/软件验证\\与确认 (V\&V)}} 
& 验证与确认概念辨析、测试战略意义、测试多维目标与测试层级、DAMP vs DRY 测试代码哲学 & 解释 & 计算机核心 & 1 \\
\cmidrule{2-5}
& 测试计划与用例生成、测试矩阵与覆盖率度量、TDD 实践、面向对象测试/Mocking/DI、代码评审、提交前后测试权衡、缺陷追踪、前沿领域验证挑战 & 解释 & 领域核心 & 3 \\
\midrule
\makecell[l]{软件工程/代码重构\\与软件演进} 
& 海勒姆定律（隐式接口定律）、向后兼容性契约本质、经典代码重构模式与 IDE 自动化重构、语义化版本控制与主干演进 & 解释 & 领域核心 & 2 \\
\midrule
\makecell[l]{软件工程/软件可靠性\\工程} 
& 可靠性数学概念、领域可靠性需求定级、缺陷注入与清除全周期、系统级冗余容错行为、现代 SRE“错误预算”与“零缺陷”对比思辨 & 解释 & 领域核心 & 2 \\
\end{longtable}
